\section{Certified fresh SLE$_6$ frontiers at stopping times}
\label{sec:frontier-certificate-bridge}

Use the comparable-cluster certificate $\mathsf C_Q$ with $d=4/3$
and the empirical decoder associated to whole-plane SLE$(8/3;2)$.
For standard SLE$_6$ in $\mathbb H$, write $K_t,g_t,W_t,\gamma_t$
and $h_t=g_t^{-1}$. Let $[\ell_t,r_t]$ be the swallowed old-real
interval and let $c_t,d_t$ be the boundary images of its endpoints.
For $t>0$, $c_t<d_t$ and $c_t\le W_t\le d_t$.

We use the following specific external inputs from
Zhan~\cite{zhan-duality-ii}: Theorem 1.1 identifies the frontier at
the swallowing time $T_p$ of fixed $p\ne0$ as a simple crosscut with
a mixed dual force-point SLE law; Lemma 6.1 gives positive probability
for either swallowing order of two real points on opposite sides;
Theorem 6.1 describes the simple fresh frontier sides at each chosen
finite stopping time. We transfer the stronger certificate below;
the dimension conclusion alone is insufficient.

\begin{lemma}[The special swallowing-time frontier is certified]
\label{lem:swallow-frontier-certified}
Almost surely every nonconstant compact interior subarc of
$\partial K_{T_p}\cap\mathbb H$ has certificate $\mathsf C_I$.
This holds simultaneously for any prescribed countable set of
nonzero real $p$.
\end{lemma}
\begin{proof}
Fix first a landing-point parameter $y$ in the duality mixture.
For $\kappa'=8/3$, the force weights are
\[
 (-4/3,0,2/3,-4/3)
 \quad\hbox{at }(0,y^a,y^b,p),\qquad
 a=\operatorname{sign}p,\quad b=-a.
 \tag{FB1}
\]
The dual trace is a simple crosscut staying in $\mathbb H$ between
its two real endpoints. Its driver has Brownian coefficient
$\sqrt{8/3}$ and drift $\sum_j\rho_j/(U-V_j)$.
A nondegenerate force point cannot be swallowed before the terminal
boundary hit: before then the hull is a simple interior slit.
The degenerate force gaps separate immediately and stay nonzero at
positive times; their weights $0,2/3$ exceed
$\kappa'/2-2=-2/3$. Consequently all force gaps are bounded away
from zero on compact time windows strictly inside the lifetime.

Restart at a positive rational time and stop before a force gap is
smaller than $1/n$, within an integer horizon. The driver drift is
bounded. Conditional bounded-drift Girsanov gives equivalence with
ordinary SLE$_{8/3}$ up to this stop. The all-affine, all-subarc
event of Theorem~\ref{thm:hereditary-cluster-certificate}, followed
by the pathwise conformal stability of
Lemma~\ref{lem:certificate-conformal-stability}, certifies the compact
future subarcs in the original dual coordinate. The inverse restart
chart is conformal near each such compact piece. Countably many
positive rational restart times, gap bounds and horizons cover all
compact interior dual subarcs; heredity handles their arbitrary
endpoints.

This is a Borel event on the unlabelled compact crosscut. Use a Borel
orientation with endpoints $u,v$; cut at the first radius-$1/n$
exit from $u$ and at the first such exit from $v$ along the reversed
arc. When $2/n<|u-v|$, the middle compact subarc $X_n$ is Borel
and is contained in the interior. Every compact interior subarc is
in $X_n$ for all large $n$. Thus the event is exactly that all these
$X_n$ have the hereditary Borel certificate. Integrating the
duality mixture is therefore legitimate. Intersecting over the
prescribed countable set of $p$ proves the assertion.
\end{proof}

\begin{theorem}[Certified fresh branches at each chosen stopping time]
\label{thm:stopping-frontier-certified}
For every chosen admissible stopping time $0<T<\infty$, almost surely
$h_T$ is injective on each of $(c_T,W_T)$ and $(W_T,d_T)$,
these images lie in $\mathbb H$, and every nonconstant compact
subarc has certificate $\mathsf C_I$. At least one of these
intervals is nonempty. No uncountable all-time intersection is used.
\end{theorem}
\begin{proof}
Take rational $a<b$ on $\{c_T<a<b<W_T\}$.
The predicate that $h_T|_{[a,b]}$ fails injectivity, fails compact
containment in $\mathbb H$, or has image outside $\mathsf C_I$ is
$\mathcal F_T$-measurable. Suppose this bad event has positive
probability. Choose a deterministic positive integer $p$ so that
its intersection with $\{r_T<p\}$ still has positive probability.
Continue by an independent centered SLE$_6$ after $T$. Conditional
on the past, put
\[
 q=g_T(p)-W_T>0,\qquad s=b-W_T<0.
 \tag{FB2}
\]
The continuation swallows $q$ before $s$ with strictly positive
conditional probability. At that time its left swallowed endpoint
is strictly greater than $s$, leaving $[a-W_T,b-W_T]$ untouched.
The original left old endpoint is unchanged.

Let $\widehat g$ be the future map and put
$F(x)=W_T+\widehat g(x-W_T)$ on the untouched interval.
It is analytic with positive real derivative there,
$F(c_T)=c_{T_p}$, and $F([a,b])$ is a compact subinterval of the
fresh crosscut interval at the original $T_p$. The inverse
composition identity gives $h_T=h_{T_p}\circ F$ on $[a,b]$.
The special-time lemma now forces injectivity, interior location,
and the certificate of its image. Thus the proposed bad event
cannot survive a conditional event of strictly positive probability.
Its intersection with the preservation event would have positive
probability, contradicting the almost-sure special-time conclusion.

Reflect for the right fresh interval, using a deterministic negative
integer $p<\ell_T$. Intersect over rational $a,b$ and the countably
many such integers. Every compact subinterval is contained in a
rational one; heredity and restriction prove the result.
The assertions apply on $\{0<T<\infty\}$ by localization.
\end{proof}

\begin{corollary}[Rough physical charts and one-sided source supply]
\label{cor:physical-stopping-frontier-certified}
Let a fixed-affine SLE$_6$ law have physical inverse chart $A\psi$,
where $\psi$ is conformal on $\mathbb H$ and $A$ is real-linear
and orientation preserving. At every chosen positive finite stopping
time before the target, at least one open prime-end boundary interval
adjacent to its source has all compact subinterval images simple,
regular, and certified by
$\mathsf C_{[(A^{-1})^{\mathsf T}A^{-1}]}$.
\end{corollary}
\begin{proof}
Each compact fresh piece in the theorem lies strictly inside the
initial $\mathbb H$. The map $\psi$ is conformal near it, so
pathwise conformal stability followed by $A$ gives its certificate.
The updated inverse chart is
$A\psi h_T(\,\cdot+W_T)$. One fresh interval adjacent to its
centered source is nonempty. No regularity of the initial physical
boundary near those compact pieces is needed. Physical trace
continuity and observable invertible capacity clocks retain their
standing admissibility meanings.
\end{proof}

\begin{lemma}[Immediate hitting of either adjacent boundary interval]
\label{lem:sle-six-immediate-side-hits}
For standard SLE$_6$, first hitting of $(0,\epsilon)$ has infimum
zero almost surely for every $\epsilon>0$, and likewise for
$(-\epsilon,0)$. Every fixed non-source real point is avoided.
\end{lemma}
\begin{proof}
The swallowing times are nondecreasing in $p>0$, finite almost
surely, and satisfy $T_p\stackrel d=p^2T_1$. Hence
$T_{2^{-n}}\downarrow0$ almost surely. Zhan's strict overshoot gives
$\gamma(T_{2^{-n}})>2^{-n}$ on the real boundary. Continuity at
zero makes these positive boundary hits tend to zero in both space
and time. Reflect for the negative side. A fixed real $p\ne0$
cannot be visited at its swallowing time, by strict overshoot, and
thereafter is target-invisible; this proves avoidance.
\end{proof}

\section{A boundary-recorded one-time change of affine metric}
\label{sec:one-switch-construction}

We give a construction on the Borel class of bounded simply connected
marked domains whose inverse disk maps extend continuously to the
closed disk. Boundaries are locally connected, but may have dense
touches and distinct prime ends with the same physical impression.
This class is closed under the restart kernels below and contains
the analytic comparison domains. Extension to arbitrary non-continuously
charted prime-end domains is not claimed.

\subsection{State coding and known-model kernels}

Use bounded univalent holomorphic inverse charts $f:\mathbb H\to D$
with continuous extension to the compactified closed half-plane.
Source and target are $0,\infty$ in the prime-end coordinate, and
the two real sides carry the two colors.
This is a Borel class in the compact-open holomorphic-function space.
After a Cayley map, continuous closed extension is equivalent to
uniform Cauchy convergence of the radial restrictions in
$C(\overline{\mathbb D})$. Boundary evaluation is a Borel limit
of interior evaluations.

Positive scalings are the only chart redundancy. An explicit Borel
slice is obtained by taking the first rational point $q(D)$ in $D$,
letting $c(f)=|f^{-1}(q(D))|$, and choosing
$\widehat f(z)=f(c(f)z)$. The injective evaluation map
$(f,z)\mapsto(f,f(z))$ and Lusin--Souslin show that inverse evaluation
and $q(D)$ are Borel. Since $c(f(d\,\cdot))=c(f)/d$, this gives
one representative per physical marked state.

Fix a nonconformal orientation-preserving real-linear $B$ and set
\[
 Q_0=I,\qquad Q_1=[B^{-\mathsf T}B^{-1}].
 \tag{OS1}
\]
Let $P^0$ be ordinary SLE$_6$ and
$P^1_D=B_*\operatorname{SLE}_6(B^{-1}D)$.
For explicit Borel dependence put $A=B^{-1}$ and let $u_f$ be the
quasiconformal self-map of $\mathbb H$ fixing $0,1,\infty$ with
coefficient
$\mu_f=\mu_A\overline{f'}/f'$. Reflection and the measurable
Riemann mapping theorem~\cite{ahlfors-bers} give existence and uniqueness. Then
\[
 F_f=A f u_f^{-1}
 \tag{OS2}
\]
is conformal onto $AD$, and a physical phase-one curve is
$f(u_f^{-1}\gamma)$ for standard SLE$_6$ $\gamma$.

The normalized $u_f$ lie in a fixed compact family of normalized
quasiconformal sphere maps. Their graph is closed: $\mu_f$ is smooth
in $\mathbb H$ and depends smoothly on compact subsets on $f$,
and its distributional equation reads
\[
 \int_{\mathbb H} u_f
 \{\overline\partial\phi-\partial(\mu_f\phi)\}\,dA=0,
 \qquad \phi\in C_c^\infty(\mathbb H).
 \tag{OS3}
\]
A countable dense family suffices, and the identity passes to uniform
limits. Uniqueness and compactness give continuity of $f\mapsto u_f$,
in particular Borelness. Its boundary extension is a homeomorphism.
Thus both known-model kernels and their restart charts are Borel.

In a baseline conformal coordinate the remaining inverse map is
$f g_t^{-1}(\,\cdot+W_t)$ and has a continuous closed extension.
Fixed affine changes preserve locally connected boundaries; the same
quasiconformal uniformization returns the restart to the stated
ordinary conformal chart class. Finite piecing is also admissible.
A quasiconformal self-homeomorphism of $\mathbb H$ preserves strict
growth of filled hulls and sends locally uniform simple approximants
to such approximants. Yuan's Theorem 1.1~\cite{yuan} gives continuous
Loewner generation in the changed coordinate; Section 2 of that
paper gives gluing through the previous inverse map with concatenated
drivers.

For completeness, full physical curves are uniform limits of embedded
arcs with converging prime-end endpoints. Approximate successively
longer pretarget prefixes by simple traces and map through the
continuous initial chart. The constructed half-plane curve tends
to infinity, since each known SLE curve does and the intervening
quasiconformal and inverse hull maps fix infinity. Physical continuity
of the initial chart at infinity makes its tail diameter tend to
zero. End each approximant at the prefix endpoint; after
reparameterization its error is bounded by the prefix error plus
the tail diameter. Endpoints approach the marked prime ends.
Truncation just inside endpoints handles coincident physical
impressions. Fixed physical approximation endpoints are not required.
The physical ordered trace determines its prime-end lift on dense
interior times by injectivity, and then everywhere by continuity;
no hidden side label is added to its filtration.

\subsection{The Borel open portion of marked old boundary}

Use the hereditary arc certificate with $d=4/3$ and put
\[
 \mathsf C=\mathsf C_{Q_1}\ \cup\
 \{X:X\text{ is a nonconstant compact simple arc on a circle
 of radius one}\}.
 \tag{OS4}
\]
The class is Borel, hereditary and translation invariant.
For the circle alternative, the relation
$X\subset\{z:|z-a|=1\}$ is closed and possible centers stay bounded
along a convergent sequence of nonempty compact $X$.
Its existence projection is therefore closed; intersect with the
Borel simple-arc class. Fresh type-zero frontier pieces fail both
alternatives: their certificate has tensor $Q_0\ne Q_1$, and smooth
finite-length arcs have zero $4/3$-dimensional tube mass.

For rational compact nondegenerate intervals
$I\subset\mathbb R\setminus\{0\}$ select those for which $f|_I$
is injective and $f(I)\in\mathsf C$, and set
\[
 O(f)=\bigcup_{I\text{ selected}}\operatorname{int}I.
 \tag{OS5}
\]
The tests and this state-dependent open set are Borel.
Heredity shows that every real compact interval in
$\mathbb R\setminus\{0\}$ gives the same union: restrict to rational
subintervals. Thus it is intrinsic under prime-end
reparameterization. Qualification is local at every interior point
of a qualifying interval: restrict its physical image into a
common neighborhood. Full colored prime-end sides are retained.

For the phase-zero curve in its initial chart let
\[
 S_f=\inf\{t\ge0:\gamma(t)\in O(f)\}.
 \tag{OS6}
\]
This is a jointly Borel stopping time, even with an unattained
infimum: exhaust selected intervals by compact inner intervals and
take the countable infimum of their closed-set hitting times.
The usual right-continuous filtration is used. Hence
\[
 \alpha(f)=P^0_f\{S_f=0\}
 \tag{OS7}
\]
is Borel and takes values only zero and one by the Brownian
initial zero-one law.

\subsection{Restart transport and absorption}

\begin{lemma}[Old-marker transport and phase-one absorption]
\label{lem:one-switch-marker-transport}
At a chosen positive finite phase-zero stopping time $T$,
\[
 O(f_T)=g_T\bigl(O(f)\cap
 ((-\infty,\ell_T)\cup(r_T,\infty))\bigr)-W_T.
 \tag{OS8}
\]
Future first visits agree with first visits to the surviving initial
open set, up to finite old/fresh seams which a future SLE$_6$
avoids. Almost surely on $\{0<S_f<\infty\}$,
$\alpha(f_{S_f})=1$. At each chosen strictly positive phase-one
stopping time before the target, its remaining ordinary-chart state
$h$ likewise has $\alpha(h)=1$ almost surely.
\end{lemma}
\begin{proof}
The remaining real boundary consists of surviving old open pieces,
one or two fresh frontier intervals, and two seams.
Every compact fresh phase-zero piece has certificate $\mathsf C_{Q_0}$.
If a selected interval intersected a fresh open interval, heredity
would certify a compact fresh subpiece, a contradiction.
A seam cannot be interior to a selected interval, since its
neighboring fresh open piece would then be intersected.
On surviving old pieces $g_T$ is an analytic real diffeomorphism,
and the physical arc image is unchanged. This proves (OS8).

Swallowed interior old boundary is target-invisible. A possible
remaining seam discrepancy is a fixed non-source prime end,
conditionally on the state, which a future standard SLE$_6$ avoids
by the preceding lemma.

At an attained activation, its real stopping tip is an old swallowed
extreme, by Zhan's Theorem 6.1. The initial open marked interval
therefore has a surviving one-sided portion adjacent to the source.
Phase-zero SLE$_6$ hits it immediately. At an unattained finite
infimum, the original phase-zero path has visits to $O(f)$
arbitrarily close from the right. Strong Markov at $S_f$, (OS8),
and seam avoidance give zero future activation infimum as well.
Thus $\alpha(f_{S_f})=1$.

At a positive phase-one stop, every compact subinterval of one
fresh source-adjacent interval has certificate $\mathsf C_{Q_1}$.
The boundary homeomorphism between its two metric charts sends
this to a full one-sided ordinary-chart interval adjacent to zero.
It is contained in $O(h)$. Immediate phase-zero boundary hitting
gives $\alpha(h)=1$.
\end{proof}

\subsection{Piecing and genuine strong DMP}

Run $P^0_f$ to $S_f$ and, if finite, independently run
$P^1_{f_{S_f}}$ thereafter. For $S_f=0$ start phase one immediately;
for $S_f=\infty$ use phase zero throughout. Borel hitting, state
and concatenation maps give a Borel physical family $Q$.

An unattained switch need not remain the same observed hitting
functional after replacing its future. Use first the expanded
piecing filtration $\mathcal G$: it reveals phase zero only up to
the switch and phase one through its current elapsed time, with
the observable exact capacity change. All completions are on the
same product probability space; adjoining its common null sets to
the phase-zero filtration does not change strong Markov.
Use initial ordinary half-plane capacity as global clock.
Before switching,
\[
 E\in\mathcal G_t\ \Longrightarrow\
 E\cap\{S>t\}\in\mathcal F^0_t.
 \tag{OS9}
\]
This persists under right-continuous augmentation: intersections
with $\{S>t+1/m\}$ are in $\mathcal F^0_{t+1/m}$; their increasing
union is the required restriction and lies in
$\bigcap_m\mathcal F^0_{t+1/m}$.

For any $\mathcal G$-stopping time $\tau$, put $\tau_0=\tau$ on
$\{\tau<S\}$ and infinity otherwise. It is a genuine phase-zero
stopping time because
\[
 \{\tau_0\le t\}=\{\tau\le t,S>t\}\ \cup\
 \bigcup_{r\in\mathbb Q\cap[0,t]}
 \bigl(\{\tau\le r,S>r\}\cap\{S\le t\}\bigr).
 \tag{OS10}
\]
The same formula with a $\mathcal G_\tau$ event inserted proves
the corresponding restriction of stopping-time information.
Phase-zero strong DMP and (OS8) therefore give conditional future
$Q_{f_\tau}$ on $\{\tau<S\}$, without a domain-Feller assumption.
At equality use $\alpha(f_S)=1$ when $S>0$ and (OS7) when $S=0$;
then $Q_{f_S}=P^1_{f_S}$. After switching, condition on switch data
and use phase-one strong DMP at the shifted stopping time with its
exact inherited clock. Absorption gives
$Q_{f_\tau}=P^1_{f_\tau}$ there too.
These cases prove the state-only strong restart identity in
$\mathcal G$.

The usual natural trace filtration is contained in $\mathcal G$.
Its stopping times are therefore $\mathcal G$-stopping times.
Taking conditional expectations, since the kernel depends only on
the observed state, proves strong DMP in the usual right-continuous
natural filtration, including time zero. Equation (OS7) excludes
a hidden random initial phase. The argument does not require
observability of the auxiliary switch from the final trace.

\subsection{Locality, translations and crossing estimates}

For full colored source neighborhoods which coincide, qualifying
old prime-end portions coincide locally by heredity. Couple the
phase-zero curves by corrected SLE$_6$ locality until exit.
Their activation rules agree before exit; if activation precedes
exit, couple phase-one continuations by the same locality.
Full boundary colors retain the same target side at contacts.
Thus corrected source locality holds also for different remote
targets. Translation covariance follows from that of the arc
class and the baselines.

Chart scaling does not affect the physical rule.
For $f_c(z)=f(cz)$,
$u_{f_c}(z)=u_f(cz)/u_f(c)$; the metric-one inverse chart changes
only by positive scaling. Also $O(f_c)=c^{-1}O(f)$.
SLE scaling and its exact clock identify the same physical prefixes.
Hence $Q$ descends to the Borel normalized physical state space.

Choose a common baseline unforced boundary-quad bound $q<1/2$
above modulus $M$. The direct SLE estimate in
\cite[Section 1.4, proof of Theorem 1.10]{ks} tends to zero with
modulus by reduction to approaching a fixed non-source boundary
point. The two fixed affine metrics change constants by bounded
quasiconformal distortion.

Split an initially unforced quad of modulus $2M$ by its middle
crosscut $C$; both marks are in the near complementary component.
Let $H_C$ be first hit, including its prime-end endpoints.
If $S\ge H_C$, a full crossing requires phase zero to cross the
near half, costing at most $q$.
If $S<H_C$, the prefix stays on the near side. The far half is
untouched, and is either swallowed or remains unforced in the
target component: the untouched connected crosscut attaches its far
pocket to only one near component and cannot reconnect two near
components separated by the prefix. Phase one's conditional crossing
bound is at most $q$. Thus the hybrid bound is $2q<1$.
Strong DMP gives the same conditional estimate at every restart.
No independence of the activation rule and the initial curve is used.

\begin{theorem}[Different affine metrics in two analytic domains]
\label{thm:one-switch-lc-counterexample}
On the explicit Borel restart-closed class of bounded locally
connected marked domains above, $Q$ has genuine strong DMP,
corrected source locality, physical translation covariance, and a
uniform unforced boundary-quad crossing bound. It has continuous
Loewner generation, strict hull growth and uniform simple
approximation with converging prime-end endpoints.
There is no deterministic affine normalizer common to all its laws.
\end{theorem}
\begin{proof}
Only the last assertion remains. In a unit disk, all old boundary
arcs satisfy the circle alternative in (OS4). Its source-adjacent
intervals are marked, so $S_f=0$ almost surely and $Q_f=P^1_f$.
In a radius-two disk, all old arcs are smooth and none lies on
a unit circle. Its marked set is empty, so $Q_f=P^0_f$.
The analytic endpoint tensor decoder of
Theorem~\ref{thm:ed-analytic} identifies respectively $Q_1$ and
$Q_0=I$ from these known-model laws. They are distinct.
A single affine normalizer would identify the same physical
tensor in both domains, a contradiction.
\end{proof}

\begin{proposition}[A continuity property respected by every fixed metric]
\label{prop:affine-sle-closed-chart-continuity}
If inverse charts $f_n,f$ converge uniformly on the closed
prime-end disk, the physical curves of each fixed affine-SLE$_6$
family admit a uniform whole-curve coupling. The new family $Q$
fails this property already along analytic disk radii tending to one.
\end{proposition}
\begin{proof}
Closed-chart uniform convergence implies compact-open convergence.
The compact closed-graph argument in (OS3) gives $u_{f_n}\to u_f$
uniformly on the sphere, and their inverses converge uniformly.
Uniform convergence of $f_n$ and uniform continuity of $f$ give
uniform physical convergence of $f_n u_{f_n}^{-1}$ to
$f u_f^{-1}$. Apply them to one common standard SLE$_6$ curve.

For $r_n\to1$, $r_n\ne1$, disks $D_{r_n}=r_n\mathbb D$ have
empty marked boundary and $Q_{D_{r_n}}=P^0_{D_{r_n}}$.
These laws converge to $P^0_{\mathbb D}$, whereas
$Q_{\mathbb D}=P^1_{\mathbb D}\ne P^0_{\mathbb D}$.
\end{proof}

\paragraph{Scope and next question.}
This is a counterexample on the explicit continuously charted,
restart-closed class, with converging approximation endpoints.
No extension to arbitrary non-continuously charted prime-end
domains is asserted. The remaining domain records the phase;
no extra color or hidden random medium is appended to the state.
A suitable domain-continuity hypothesis excludes this example,
but its sufficiency for classification is open. Earlier rigidity
results conditional on a fixed-tensor mixture representation retain
that hypothesis: a general law here has a geometric change of metric.
The bare axioms on this class cannot force the common signed
regularity needed for the earlier classification theorems.
